\section{Synchronization of tensor parameters}\label{sec:sync}

Theorem~\ref{thm:extension} tests whether one price extension works for
an entire family. Its tests allow one to choose any tensors in the family,
multiply them and close their extra legs jointly. For a family depending
on a marked point $s$, two factors in such a test can therefore have
different marked points. The sphere calculation in Section~\ref{sec:sphere}
will give estimates when all factors use the same marked point.
We show that under a strict saving for square multiplication this
diagonal test already controls every choice of the points.

Fix an algebraically closed field $F$, an ordinary price $H$ of scale
$D>0$ on base sites $S$, and a finite set $L$ of additional sites,
called openings. Recall from Equation~\eqref{eq:prices-beta} that
$\beta_H$ measures the largest price per copy obtained by taking a
tensor power and applying one joint functional at each opening.
For a polynomial tensor family $T(\mathbf s)$, $\mathbf s=(s_1,\ldots,s_d)$,
write
\[
 \beta_H^{\rm gen}(T(\mathbf s))
    =\sup_{\mathbf c\in F^d}\beta_H(T(\mathbf c)).
\]
This is the assisted cost at independent indeterminates for the
generic extension of $H$ in Section~\ref{sec:valued}. Indeed, fix a
power and a closure with rational function coefficients. One can
specialize its entries together with the tensor on a common open
set where the resulting price is arbitrarily close to its generic
value. Conversely, a fixed closure over $F$ witnessing a strict lower
bound at one point witnesses the same bound on a nonempty open set,
by closedness of price sublevels. In particular the supremum is unchanged
on restricting to any dense open parameter domain.

\begin{theorem}[Synchronization]\label{thm:sync}
Suppose $\omega_F<9/4$. For any polynomial tensor families
$U_1,\ldots,U_a$ on $F^d$ with common base and openings,
\[
 \beta_H^{\rm gen}\!\left(\bigotimes_{i=1}^a U_i(\mathbf s_i)\right)
 =
 \beta_H^{\rm gen}\!\left(\bigotimes_{i=1}^a U_i(\mathbf s)\right).
\]
On the left the $a$ parameter vectors vary independently. On the right
they are equal. Both products group their legs at each opening, so a
closure on either side may couple all the factors.
\end{theorem}

\begin{proof}
Choosing equal parameters gives the inequality from right to left.
For the reverse inequality we first take copies of one univariate
family $U(s)$. An identically zero factor makes the assertion immediate,
so we may assume that the families under consideration are nonzero.
Write $B=\beta_H^{\rm gen}(U(s))$. The diagonal cost of $a$ copies is
$aB$ by homogeneity under tensor powers. We first control different
parameters lying in a disk above the same residue $s$, then recover
independent parameters by letting the valuation depth of that disk
tend to zero.

\smallskip\noindent\emph{Pinning above one residue.}
Take an uncountable set $\mathcal U$ of independent indeterminates,
and let $s$ be a further indeterminate. Put
\[
 E=F(\mathcal U,s),\qquad
 K_0=\operatorname{Frac}(F(\mathcal U)[t^{\mathbb Q}]),\qquad
 J=\{z\in K_0:\nu(z)>0\}.
\]
The Gauss lift of the generic extension of $H$ to $E$ is a valued price
$p$ on $K_0(s)$ (Theorem~\ref{thm:gauss}). For $z\in J$, translation
by $z$ preserves the unit disk. Apply Lemma~\ref{lem:star} to this
translation and its inverse, taking $s$ as the fresh residue variable.
It gives
\begin{equation}\label{eq:sync-disk-single}
     \beta_p(U(s+z))=\beta_p(U(s))=B .
\end{equation}
For the second equality the Gauss lift agrees with the ordinary
generic price on coefficients independent of $t$. As in the proof of
Lemma~\ref{lem:star}, it suffices to use closing functionals with
constant entries in $F$.

Fix $M\ge1$ and an anchor in $S$. For each $z\in J$, let $V_z$ be the
product of $U(s+z)^{\otimes M}$ with the pin penalties of
Lemma~\ref{lem:pinning}, one for each opening. By translation invariance
its minimum $b=\beta_p(V_z)$ is the same for every $z$. Give each
objective its own copy of $L$. Lemma~\ref{lem:pinning} realizes its minimizing
sites $x_{z,\ell}$, $\ell\in L$, in one complete $\operatorname{CAT}(0)$
spectral space over $F$. For any finite list of objectives, a valued
extension simultaneously realizes their minima.

\smallskip\noindent\emph{The minimizing sites cannot move.}
For each $\ell$ we check that a nonconstant map $z\mapsto x_{z,\ell}$
would satisfy the hypotheses of Theorem~\ref{thm:motion}. The bridge
from these ordinary sites to the valued disk is the pinned evaluation
formula. For every nonzero test tensor $T$ over $F$ placed at sites
from a finite list $z_1,\ldots,z_m$, its price is
\begin{equation}\label{eq:sync-pinned-prices}
 \lim_{n\to\infty}
 \left[
 \beta_p\!\left(
  (V_{z_1}\otimes\cdots\otimes V_{z_m})^{\otimes n}\otimes T
 \right)-nmb
 \right].
\end{equation}
The $m$ copies of $L$ here are distinct. If two parameters coincide,
the corresponding minimizing sites are at distance zero, by uniqueness
of pinning.

Permutation of the indeterminates in $\mathcal U$ preserves the prices
in Equation~\eqref{eq:sync-pinned-prices} after relabeling the $z_j$.
Translating every $z_j$ by $w\in J$ replaces $s$ by $s+w$ in a
polynomial family over $K_0$, so the disk comparison preserves each
finite-$n$ expression as well. For an affine test $z_j=a_j+b_jc$ in
Theorem~\ref{thm:motion}, the tensor inside $\beta_p$ is polynomial in
$c$. Lemma~\ref{lem:star} therefore gives specialization and the
weighted submean inequality at the disk and its inner and outer
ends. Subtracting the same constant $nmb$ at every end preserves these
inequalities, and taking the limit transfers them to the pinned prices.

It remains to check continuity at zero. At any one placement of the
openings, $U(s+z)$ and $U(s)$ both have an entry of valuation zero,
because their common reduction $U(s)$ is nonzero. Each tensor has
nonnegative price there. The entries of their difference have
valuation at least $\nu(z)$. Applying the entry bound
\eqref{eq:prices-range} to the difference and the sum inequality in
both directions shows that the prices of the two tensors differ by
at most
\[
      \eta(z)=D\log\!\left(1+\exp((C-\nu(z))/D)\right)
\]
for a constant $C$ depending only on the format and scale. Evaluate
$V_0$ at the sites minimizing $V_z$. Its price is at most
$b+M\eta(z)$ since the pin penalties are identical at that placement.
In a common minimizing extension for the two lists, the midpoint
estimate \eqref{eq:prices-midpoint} applied to $V_0$ gives
\[
  2c_D\sum_{\ell\in L}d(x_{z,\ell},x_{0,\ell})^2 \le M\eta(z).
\]
The right side tends to zero as $\nu(z)\to\infty$ at fixed $M$.
Distances are bounded by the scale $D$. A nonconstant family of
pinned sites would now force exact matrix multiplication rank at
least $n^{9/4}$ for every $n$, by Theorem~\ref{thm:motion}.
Lemma~\ref{lem:bilinearize} would then give
$\omega_F\ge9/4$. Thus for every $\ell$, all $x_{z,\ell}$ coincide.

\smallskip\noindent\emph{Removing the penalty and recovering independent parameters.}
The sum of the penalties is between zero and a constant $C_0$
depending only on $D$ and $|L|$. Using a placement minimizing
$U(s+z)$ as a competitor gives $b\le MB+C_0$ by
Equation~\eqref{eq:sync-disk-single}. Hence at its pinned sites
$U(s+z)$ has price at most $B+C_0/M$. For any $z_1,\ldots,z_a\in J$,
take a common valued extension minimizing their pin objectives.
The corresponding openings have distance zero. Transport therefore
identifies their prices also on tensors with coefficients in $K_0(s)$.
Product additivity at these common openings gives
\[
  \beta_p\!\left(\bigotimes_{i=1}^aU(s+z_i)\right)
       \le a(B+C_0/M).
\]
Let $M\to\infty$ to obtain $aB$ as an upper bound throughout the disk.

Take independent $u_1,\ldots,u_a\in\mathcal U$. Fix a power $h$ and
a closure over $F$, and apply this closure to
$(\bigotimes_i U(s+v u_i))^{\otimes h}$, obtaining a polynomial
tensor $P(v)$ over $E$ on the base sites. For rational $r>0$ the
parameters $s+t^r u_i$ lie above $s$, so $p(P(t^r))\le haB$.
Write $H_E$ for the ordinary generic extension of $H$ to $E$.
If $P(1)\ne0$, the Gauss formula gives
\[
                 H_E(P(1))\le p(P(t^r))+r\deg P .
\]
Indeed, the coefficients of $P(v)^{\otimes n}$ sum to $P(1)^{\otimes n}$.
There are at most $n\deg P+1$ coefficients, and the Gauss penalty
for any of them is at most $rn\deg P$. The loss from adding them is
at most $D\log(n\deg P+1)$. Divide by $n$ and pass to the limit using
Equation~\eqref{eq:gauss-price}. Letting $r\downarrow0$ gives
$H_E(P(1))\le haB$. The substitution $u_i\mapsto s+u_i$ is an invertible change of
variables over $F(s)$, so $P(1)$ tests the independent generic
parameters.
Taking all powers and constant closures proves the identical-family case.

\smallskip\noindent\emph{Different families and several coordinates.}
For arbitrary univariate families $U_i$, apply the result just proved to
$U(s)=\bigotimes_iU_i(s)$. With independent $s_1,\ldots,s_a$, regroup
$\bigotimes_jU(s_j)$ as the product of $a$ tensors
$\bigotimes_i U_i(s_{i+j})$, reading indices cyclically. Each of these
has the same assisted cost at the generic price on $F(s_1,\ldots,s_a)$
because that price is invariant under permutations of the
indeterminates. Superadditivity of assisted cost on products follows by taking
separate closures of the factors on a common multiple of their
witnessing powers. Thus $a$ times the independent cost for the $U_i$
is at most the independent cost for $a$ copies of $U$, which we have
bounded by $a\beta_H^{\rm gen}(U(s))$.

For several coordinates choose any $\mathbf c_1,\ldots,\mathbf c_a\in F^d$.
Coordinatewise interpolation gives a single polynomial curve
$\gamma:\mathbb A^1\to\mathbb A^d$ through these points at distinct
scalar arguments. The univariate result applied to $U_i\circ\gamma$
bounds the assisted cost of this tuple by the common-parameter cost
along $\gamma$, hence by the common-parameter cost on $F^d$.
Taking the supremum over tuples completes the proof.
\end{proof}

\begin{corollary}[Simultaneous pricing]\label{cor:sharedprice}
Under the hypotheses of Theorem~\ref{thm:sync}, let
$\{U_i(\mathbf s)\}_{i\in I}$ be a collection of polynomial families
with common openings. For every finite list $i_1,\ldots,i_m$,
allowing repetitions, suppose every closure of
$\bigotimes_j U_{i_j}(\mathbf s)$ at a shared generic parameter has
price at most $\sum_j h_{i_j}$. Then one ordinary extension
$\widehat H$ of $H$, of the same scale, satisfies
$\widehat H(U_i(\mathbf c))\le h_i$ for every $i$ and every
$\mathbf c\in F^d$. The conclusion also holds for rational families
on a common nonempty open domain of regularity. In particular, for
one such family $U$ on its domain $\Omega$,
\[
 \min_{\widehat H\supset H}\ \sup_{\mathbf c\in\Omega}\widehat H(U(\mathbf c))
    =\sup_{\mathbf c\in\Omega}\beta_H(U(\mathbf c)).
\]
\end{corollary}

\begin{proof}
Repeating each finite list gives the assumed bound on closed powers
and hence on its synchronized assisted cost. By Theorem~\ref{thm:sync}
this bound also holds at independently chosen parameters.
Theorem~\ref{thm:extension} now yields the single price extension.

For a finite list of rational families, clear scalar denominators on a
dense open set where they are nonzero. An ordinary price is unchanged
by nonzero scalar multiples, so the polynomial result gives the same
product tests there. If such a test failed at a regular tuple, a
fixed power and closure would witness a strict violation. Closedness
of price sublevels makes this violation persist on an open neighborhood
of the tuple, which meets the dense open set. The tests now hold
throughout the domain, as required by Theorem~\ref{thm:extension}.

For the minimax identity use the supremum of the pointwise assisted
costs as the bound for the single family. If the family is
nonzero this is finite by the fixed-format bound. Products at a
shared parameter satisfy the tests by homogeneity of $\beta_H$ on
powers, giving an extension with the stated upper bound. Every
extension majorizes $\beta_H$ pointwise, giving the reverse inequality.
The zero family has both sides equal to $-\infty$.
\end{proof}

For the sphere tensor, $s$ specifies the additional marked point and
the family index specifies its energies. Corollary~\ref{cor:sharedprice}
places its leg once for all regular positions and all energies. Each
collision with an original marked point can then measure a distance
from this same site.
